% Compiles with pdfLaTeX. No CJK packages needed.
% Rewritten 2026-08-17 against the part1c / part2 / part3 archives.
% Every number here was read from those archives during the rewrite. Nothing was carried
% over from the previous notes without re-verification; the places where the previous notes
% were wrong are marked explicitly so the two versions can be diffed.
\documentclass[10pt,a4paper]{article}

\usepackage[T1]{fontenc}
\usepackage[utf8]{inputenc}
\usepackage{lmodern}
\usepackage[margin=0.8in]{geometry}
\usepackage{amsmath,amssymb}
\usepackage{booktabs,array,multirow,graphicx}
\usepackage{xcolor}
\usepackage{enumitem}
\usepackage[colorlinks=true,linkcolor=blue!55!black,urlcolor=blue!55!black]{hyperref}
\usepackage{titlesec}

\setlist{nosep,leftmargin=1.7em}
\titlespacing*{\section}{0pt}{1.2em}{0.5em}
\titlespacing*{\subsection}{0pt}{0.85em}{0.35em}
\renewcommand{\arraystretch}{1.15}
\newcommand{\ours}{\textsuperscript{$\dagger$}}
\newcommand{\fix}[1]{\textcolor{red!65!black}{#1}}

\title{\textbf{EfficientEval Experiment Notes, v4}\\[3pt]
\large One routing method, two evaluation protocols, rewritten against the frozen archives}
\author{}\date{2026-08-17}

\begin{document}
\maketitle
\vspace{-2.2em}
\hrule

\section*{Status of this document}

This replaces the previous notes. The previous version described a configuration that no longer
exists: nine cheap features for Protocol A and five for Protocol B, a 6{,}850-row Protocol~A
frame, $p\ge0.5$ as the decision rule, and result tables produced under all of
those. The routing algorithm is unchanged. Everything downstream of it was re-selected, re-run,
or re-derived, so no table from the previous version survives.

Provenance for every number below:
\begin{itemize}
\item main results --- \texttt{artifacts/part1c\_main\_full\_v1}
\item core ablation --- \texttt{artifacts/part2\_ablation\_v1}
\item extended ablation --- \texttt{artifacts/part3\_extended\_v1}
\item pool provenance --- \texttt{runs/pool\_rescreen\_v1}
\item verifier cost structure --- \texttt{DELIVERABLE/05\_verifier\_code/COMPLEXITY.md}
\end{itemize}

Two result groups were still computing when this was written and are marked
\textbf{[pending]}: the Protocol~A replication of the extended ablation, and the cascade,
second-call and learner families. Their Protocol~B counterparts are complete and reported.

Corrections to the previous notes are marked \fix{in this colour}.

%==================================================================
\section{Task and data}

An instance is a source document $D_i$, a candidate summary $S_i$, and a human factuality label
$y_i\in\{0,1\}$. The router sees only features of $(D_i,S_i)$, chooses one of $k=3$ verifiers,
calls it, and emits a calibrated probability plus a binary decision.

\subsection{The two protocols}

\begin{center}\small
\begin{tabular}{p{3.2cm}p{5.5cm}p{5.5cm}}
\toprule
 & \textbf{Protocol A: pooled rotation} & \textbf{Protocol B: official TRAIN$\to$TEST} \\
\midrule
Frame & \fix{8{,}512 rows / 1{,}535 groups} (TRAIN $\cup$ TEST) &
  TRAIN 5{,}276 / 890; TEST 3{,}236 / 645 \\
Positive rate & \fix{0.5721} & TRAIN 0.6298; TEST 0.4781 \\
Source words & median 453, p90 1{,}218, max 8{,}847 &
  TRAIN 462 / 1{,}402 / 8{,}635; TEST 439 / 991 / 8{,}847 \\
Summary words & median 71, p90 151, \fix{max 2{,}099} &
  TRAIN 82 / 159 / 1{,}428; TEST 57 / 124 / 2{,}099 \\
Fitting & 8 folds fit, 1 validates & each seed splits TRAIN by group into fit and validation \\
Evaluation & the remaining fold; ten rotations concatenate out-of-fold & the frozen model on TEST \\
Repetitions & 10 rotations $\times$ 10 seeds & 10 seeds \\
\bottomrule
\end{tabular}
\end{center}

\fix{Three corrections.} The previous notes gave Protocol~A as 6{,}850 rows / 1{,}241 groups with
a 52.3\% positive rate and a maximum summary length of 1{,}444 words, and stated that the
longest summary in Protocol~B's TEST side (2{,}155 tokens) exceeded anything Protocol~A
contained. Protocol~A is the union of TRAIN and TEST, so it contains that summary; the two
frames cannot differ in their maximum. The current figures are above.

\fix{A fourth correction.} The previous notes described TEST access as ``labels are read once''.
That is not the contract. The contract is that TEST never participates in any choice --- not in
features, pool, target, learner, hyperparameters, the beta rule, calibration, or thresholds.
How many times a frozen evaluation script reads the labels is not a constraint and was never
enforced as one.

Both protocols enforce group separation on \texttt{content\_doc\_key}:
\[
G_{\mathrm{fit}}\cap G_{\mathrm{val}}=\varnothing,\quad
G_{\mathrm{fit}}\cap G_{\mathrm{eval}}=\varnothing,\quad
G_{\mathrm{val}}\cap G_{\mathrm{eval}}=\varnothing,
\]
asserted at run time rather than by convention, because one source document carries several
candidate summaries and row-level splitting would leak it. TRAIN and TEST share no document.
Five RAGTruth TRAIN rows carrying a CNN/DailyMail article that also appears in CogenSumm TEST are
excluded by a content-hash rule declared before either file existed.

\subsection{Protocol A is not a held-out confirmation}

Protocol~A's configuration was selected on the TRAIN subset of its own pooled frame, so A is a
cross-validated estimate. Protocol~B is the confirmatory result. This is stated here because the
distinction is easy to lose: both protocols report the same metrics on frames that overlap.

%==================================================================
\section{Method}

\subsection{Offline training, then one call at inference}

\begin{center}\small
\begin{tabular}{llp{7.6cm}}
\toprule
Step & Partition & What it does \\
\midrule
T1 & --- & read frozen scores $s_{ij}$, availability, and measured latency $c_j$ \\
T2 & fit & fit Platt $(a_j,b_j)$ per verifier, giving calibrated $C_{ij}$ \\
T3 & fit & build the supervision target $T_{ij}$ from $C_{ij}$ and $y_i$ \\
T4 & fit & fit three random-forest regression heads on the six features \\
T5 & validation & sweep the beta grid; select $\beta^{*}$ \\
T6 & validation & fit the stage-2 isotonic layer under $\beta^{*}$ \\
\fix{T7} & \fix{validation} & \fix{select the decision threshold} \\
\bottomrule
\end{tabular}
\end{center}

\fix{Correction.} The previous notes listed four training artifacts
$\{(a_j,b_j)\},\{f_j\},\beta^{*},\mathrm{Iso}$ and gave the decision rule as $p\ge0.5$. The
deployed object also contains a threshold selected on validation, and the primary rule is the
validation MCC-optimal cut, not $0.5$. Section~\ref{sec:thr} reports all four threshold rules
that were computed, and one of them is $0.5$.

At inference the router computes the six features, evaluates three heads, forms
$U_j=\hat T_j\exp(-\beta^{*}c_j/c_{\max})$, calls $\arg\max_j U_j$ and nothing else, applies that
verifier's Platt map, applies the isotonic layer, and thresholds. Exactly one verifier call per
instance; \texttt{mean\_calls} is $1.000$ for every primary row in every table.

\subsection{Unified cheap features}

\fix{The largest change.} The previous notes described protocol-specific feature sets, nine for A
and five for B, selected separately. There is now one set of six, selected on TRAIN alone by
jointly optimising A-style and B-style pooled AUROC:

\begin{center}\small
\begin{tabular}{p{4.6cm}p{10.4cm}}
\toprule
Feature & Definition \\
\midrule
\texttt{structured\_source\_line\_ratio} & structured lines as a fraction of non-empty source
  lines; a line is structured if it matches a bullet or numbering pattern, or has $\le8$ tokens
  and either ends in a colon or is all-uppercase \\
\texttt{bm25\_mean3} & mean BM25 score of the summary against its three best-matching source
  passages \\
\texttt{entity\_coverage} & fraction of summary entities present in the source; $1$ when the
  summary has no entities \\
\texttt{entity\_value\_colocation} & rate at which an entity and a numeric value that co-occur in
  a summary sentence also co-occur in some source sentence \\
\texttt{year\_count} & count of four-digit year mentions in the summary \\
\texttt{conflicting\_value\_rate} & rate of summary numeric values that appear in the source with
  a different associated entity \\
\bottomrule
\end{tabular}
\end{center}

Section~\ref{sec:unify} reports the head-to-head against both superseded sets. The unified set is
separably better than each of them, so unification was not a simplification bought with quality.

\fix{One caveat the previous notes did not carry.} Feature selection scored candidates on
pre-stage-2 routed AUROC. It did not fit an isotonic layer per candidate, so the search did not
execute the full two-stage pipeline it was selecting for.

\subsection{Supervision target}

With $Q_{ij}$ the probability mass verifier $j$ placed on the true label,
\[
Q_{ij}=\begin{cases}C_{ij}&y_i=1\\ 1-C_{ij}&y_i=0\end{cases}
\qquad
T_{ij}=1+\bigl(Q_{ij}-\max_k Q_{ik}\bigr).
\]
The best verifier on a row gets exactly $1$; the others decay by their regret. Absolute quality
conflates instance difficulty with verifier differences, and the row-wise maximum cancels the
shared difficulty term while preserving both the ordering and the magnitude of the gaps.

The five alternatives measured in Section~\ref{sec:abl} are \texttt{prob\_quality} $=Q_{ij}$;
\texttt{brier} $=1-(C_{ij}-y_i)^2$; \texttt{pairwise\_rank}, the fraction of other actions
verifier $j$ beats on that row; \texttt{is\_best}, the indicator of attaining the row maximum;
and \texttt{correctness}, the indicator that the hard decision matches the label.

\fix{Note on \texttt{correctness}.} It needs per-verifier hard-decision correctness, which the
frozen matrix does not carry as a column. It is derived from the recorded \texttt{decision\_\_}
columns against the label. It is the only ablation arm whose inputs are computed rather than read.

\subsection{Random-forest heads and how their hyperparameters were chosen}

One independent regression head per verifier, $f_j:\mathbb{R}^{6}\to[0,1]$, each fitted only on
the fit rows where that verifier is available, predictions clipped to $[0,1]$, split criterion
squared error.

\begin{center}\small
\begin{tabular}{lcc}
\toprule
 & Protocol A & Protocol B \\
\midrule
\texttt{n\_estimators} & 800 & 200 \\
\texttt{min\_samples\_leaf} & 5 & 10 \\
\texttt{max\_features} & 0.5 & 0.5 \\
\texttt{max\_depth} & 6 & 6 \\
\bottomrule
\end{tabular}
\end{center}

\fix{Correction, and the reason this section exists.} The previous notes stated that each protocol
chose hyperparameters ``on validation by minimising head loss over the declared 16-point search''.
No such search had been run for the feature set in use: Protocol~A's values were a hard-coded
default inherited from an earlier codebase, and Protocol~B's came from a search over a different,
five-feature set. Both have now actually been selected, on TRAIN only, by minimum validation head
loss over sixteen randomly initialised candidates drawn from
$\texttt{n\_estimators}\in\{200,400,800\}$, $\texttt{min\_samples\_leaf}\in\{1,2,5,10\}$,
$\texttt{max\_features}\in\{\texttt{sqrt},\texttt{log2},0.5\}$,
$\texttt{max\_depth}\in\{6,10,12,16,24\}$, with the two previous values forced in as reference
points. The full ranking is archived.

\fix{The cost of that rule is recorded rather than hidden.} Head loss and routed AUROC are
positively rank-correlated across the candidate set, so the minimum-loss region is the
shallow-forest region: both protocols select \texttt{max\_depth}$=6$. Selecting on validation
AUROC instead would have gained a small amount of validation AUROC. Head loss was kept because it
is disjoint from the metric the paper reports, and the forgone amount is archived alongside.

Section~\ref{sec:conv} shows the forests are past convergence at these sizes.

\subsection{Cost-aware selection}

$c_j$ is verifier $j$'s mean latency on the fit partition and $c_{\max}$ the largest of them.
On TRAIN the three are \fix{$49.4$, $81.1$ and $205.3$\,ms}, giving:

\begin{center}\small
\begin{tabular}{lrrrl}
\toprule
$\beta$ & FactCC & Lettuce-v2 & Granite-3.1-2b & Interpretation \\
\midrule
0.0 & 1.000 & 1.000 & 1.000 & cost ignored \\
0.1 & 0.976 & 0.961 & 0.905 & selected in both protocols \\
0.2 & 0.953 & 0.924 & 0.819 & Granite must predict ${\sim}16\%$ better to win \\
0.4 & 0.908 & 0.854 & 0.670 & ${\sim}36\%$ better \\
\bottomrule
\end{tabular}
\end{center}

$\beta^{*}$ is the lowest-latency value whose validation AUROC is within $\varepsilon=0.005$ of
the best over the grid $\{0,0.1,0.2,0.4,0.6,0.8,1.2,1.6,2.4,3.2\}$. Mode $0.1$ in both
protocols: 8 of 10 seeds in B, 51 of 100 rotations in A with $0.2$ taking the other 49.

\fix{Do not confuse two statistics.} The previous notes reported a Protocol~A beta selection rate
of 72\%. That was measured during the TRAIN-only selection stage, over a different set of fits,
and is not comparable to the 51\% over the 100 final rotations.

\subsection{Two calibration stages, applied to every system}

Stage 1 is a per-verifier Platt map fitted on the fit partition. Stage 2 is one isotonic layer
fitted on the validation partition, on the router's own output.

\fix{The correction that matters most for reading the tables.} In the previous work only the
router received stage 2; the fifteen fixed verifiers received stage 1 alone. That made ECE, MCE,
Brier and every threshold-dependent column incomparable between the router and its baselines, and
it produced the claim that the router attains the lowest calibration error in the table. Under
symmetric calibration that claim does not survive: in Protocol~A three fixed verifiers have lower
ECE than the router, and in Protocol~B one does. All sixteen systems now go through the identical
two-stage pipeline.

Isotonic regression is monotone but \emph{not} order-preserving. PAVA pools the score regions
where the empirical label rate runs backwards, and every pair inside a pooled block becomes a tie
worth exactly $0.5$, so a block that was contributing less than $0.5$ gains. Stage 2 therefore
moves AUROC in either direction, and moves it most for the weakest verifiers.
Section~\ref{sec:abl} quantifies this for every system.

\subsection{Latency accounting}
\label{sec:lat}

\fix{Corrected and now complete.} Reported latency is end to end: base feature extraction, the
second feature extractor, random-forest head inference, the routing arithmetic, both calibration
stages, and the verifier call. The previous accounting charged only the verifier call plus base
feature extraction, understating the router by 23\% in Protocol~B and 28\% in Protocol~A.

A fixed verifier computes no features and runs no heads, so its end-to-end cost is its call plus
the calibration applied to it. That asymmetry is correct accounting, not a handicap, and it is
what the tables report. Measured components, per row:

\begin{center}\small
\begin{tabular}{lrr}
\toprule
Component & Protocol B & Protocol A \\
\midrule
base feature extraction & 7.665 & 9.334 \\
second feature extractor & 19.628 & 23.888 \\
head inference (three forests) & 0.034 & 0.483 \\
routing arithmetic & 0.00005 & 0.00005 \\
calibration application & ${\approx}0$ & ${\approx}0$ \\
selected verifier call & 77.78 & 78.55 \\
\midrule
total & \textbf{106.14} & \textbf{112.54} \\
\bottomrule
\end{tabular}
\end{center}

The breakdown reconciles with the reported mean to $10^{-13}$\,ms, asserted at run time.

\fix{One measurement caveat that the previous notes did not contain.} Granite-3.1-2b's recorded
latency differs by 1.5 to 1.7$\times$ between two scoring campaigns covering different corpora,
and it carries 30.9\% of Protocol~B traffic. A matched-token-bin comparison on the one corpus
family present in both campaigns gives a ratio of 1.003, which rules out input length and split
composition as the explanation. The tables report the as-measured figure, which is the
conservative end; normalising every row to the slower campaign widens the router's advantage over
the strongest pool member from 1.24$\times$ to 1.51$\times$. No system dominates the router under
either accounting.

%==================================================================
\section{Verifier cost structure}

The symbolic derivations in the previous notes --- forward-pass counts, per-pass cost, the
parameter-substituted expansion, and the model-structure table --- depend only on model
architecture and the executed code path. That code path was re-audited line by line for this
rewrite and the derivations stand, with two corrections. They are not reproduced here; they live
with the code, in \texttt{DELIVERABLE/05\_verifier\_code/COMPLEXITY.md}, together with the
dispatch table and the file map.

\fix{Correction 1.} HHEM and AlignScore are not \texttt{PairWindowScorer}s. They are
\texttt{MetadataScorer}s wrapping scorers built by \texttt{build\_primary\_scorer}, and HHEM's
windowing is a separate implementation, \texttt{build\_hhem\_token\_windows}, in a different
file. The previous notes described HHEM's windowing correctly but filed it under the wrong
wrapper, so a reader following the dispatch would not find it.

\fix{Correction 2.} The closed form
$C=\lceil (m-o)/(T_{\max}-E(b)-o)\rceil$ is a lower bound, not an exact count. After slicing
\texttt{budget} token ids, the packer re-encodes the \emph{decoded} chunk as a pair and shrinks
the window while that length exceeds $T_{\max}$. Decoding and re-encoding is not
length-preserving, so the true count can exceed the closed form. The recorded
\texttt{source\_window\_count} is the ground truth.

Two facts from that document are load-bearing for any complexity claim. First, three columns
describe forward work and they disagree by design: \texttt{model\_forward\_calls} is always $1$
for both MiniCheck variants because their windows are packed into one batch, and
\texttt{forward\_items} degenerates to $1$ for Lettuce-v2 because its wrapped scorer never writes
\texttt{n\_chunks}. The correct column depends on the verifier group. Second, every scorer
asserts batch size one, so per-instance latency is measurable but nothing here is
batch-throughput-optimal.

All \emph{measured} quantities in the previous notes' cost tables --- forward-count quantiles,
median latencies, generation-length distributions, and the latency-versus-length correlation
table --- were taken from the superseded 6{,}850-row matrix and do not transfer. They can be
recomputed from the frozen scoring parquets without re-running any verifier.

%==================================================================
\section{Metrics}

AUROC is primary and the only basis for ranking. AUPRC is reported with unsupported as the
positive class. ECE and MCE use 15 equal-width bins. AURC is the area under the risk--coverage
curve, with coverage ordered by $|p-0.5|$. Latency is reported as mean, p50, p95 and p99. Macro
AUROC, the unweighted mean over the four corpora, is a diagnostic column throughout and never a
ranking key.

\fix{The composite metric is gone.} An earlier quality--cost composite used fixed normalisation
anchors that the observed quality range later exceeded, which broke the geometry it was built on.
It is removed from every table, ranking and conclusion. Quality and cost are reported separately.

\fix{Multiplicity is now reported.} Each main table makes fifteen comparisons against the router
and each ablation table makes twenty-five or more. Every paired interval is reported at the
nominal 95\% level and at the Bonferroni-corrected level for its family, both from one cluster
bootstrap resampled by \texttt{content\_doc\_key} with 2{,}000 draws. Claims cite the corrected
level.

%==================================================================
\section{Main results}

\subsection{Protocol A}

\begin{center}\scriptsize
\setlength{\tabcolsep}{2.6pt}
\resizebox{\textwidth}{!}{%
\begin{tabular}{lrrrrrrrrrrrrr}
\toprule
System & AUROC$\uparrow$ & SD$\downarrow$ & AUPRC-U$\uparrow$ & Acc$\uparrow$ & BAcc$\uparrow$ & MCC$\uparrow$ & F1-U$\uparrow$ & Rec-U$\uparrow$ & ECE$\downarrow$ & Brier$\downarrow$ & Mean$\downarrow$ & P95$\downarrow$ & P99$\downarrow$ \\
\midrule
\textbf{Ours} & \textbf{.79330} & .00173 & \textbf{.78383} & .73793 & .71951 & .45888 & \textbf{.65862} & \textbf{.59176} & .02161 & \textbf{.17262} & \textbf{112.54} & 313.00 & 1014.45 \\
AlignScore & .78113 & .00131 & .76517 & .71236 & .68446 & .40741 & .59303 & .49110 & .02154 & .18128 & 677.29 & 2195.91 & 5275.71 \\
Qwen30-fast & .76893 & .00319 & .77657 & .74934 & .71863 & .49619 & .63315 & .50574 & .01440 & .17332 & 273.35 & 506.33 & 1334.30 \\
Qwen30-judge & .76390 & .00206 & .75124 & \textbf{.75145} & .71899 & \textbf{.50524} & .62963 & .49401 & .02029 & .17310 & 642.74 & 1394.52 & 1777.29 \\
FactCG & .75500 & .00224 & .75869 & .73033 & .69847 & .45352 & .60232 & .47757 & .02079 & .18411 & 429.94 & 1562.36 & 4053.75 \\
Granite-3.1-2b\ours & .74024 & .00129 & .72905 & .70638 & .67233 & .39997 & .55968 & .43630 & .02113 & .19460 & 177.34 & 325.18 & 912.45 \\
WeCheck & .73337 & .00189 & .72338 & .69851 & .66035 & .38843 & .52868 & .39585 & .02147 & .19581 & 359.32 & 1296.06 & 3240.79 \\
MiniCheck-FT5 & .70117 & .00182 & .66240 & .68440 & .65724 & .34345 & .55968 & .46897 & .02393 & .21020 & 451.01 & 1750.79 & 4330.68 \\
FactKB & .67432 & .00255 & .69030 & .69974 & .65272 & .42093 & .48222 & .32680 & .02501 & .19999 & 44.46 & 150.52 & 385.34 \\
Granite-3.2-8b & .67028 & .00547 & .68705 & .70320 & .65685 & .42766 & .49176 & .33561 & \textbf{.01762} & .19798 & 1011.32 & 1649.83 & 3125.17 \\
Granite-3.2-3b & .66372 & .00202 & .64482 & .66350 & .61813 & .31240 & .43568 & .30365 & .02103 & .21520 & 243.92 & 465.30 & 785.93 \\
MiniCheck-DBTA & .63642 & .00191 & .59806 & .61658 & .56662 & .19325 & .32609 & .22037 & .02407 & .22554 & 473.33 & 2046.78 & 4200.11 \\
HHEM & .62453 & .00225 & .60837 & .63696 & .58399 & .25540 & .33807 & .21686 & .01739 & .22438 & 149.64 & 542.78 & 1676.51 \\
Granite-4.1-3b & .61887 & .00177 & .54875 & .58610 & .57546 & .15185 & .50862 & .50167 & .02195 & .23446 & 448.81 & 657.21 & 1449.36 \\
Lettuce-v2\ours & .61216 & .00488 & .55749 & .65578 & .61229 & .28777 & .43590 & .31085 & .01403 & .22607 & 73.68 & 233.09 & 722.36 \\
FactCC\ours & .53237 & .00287 & .45330 & .51523 & .51612 & .03225 & .47738 & .52230 & .01972 & .24544 & 44.86 & 149.36 & 374.19 \\
\bottomrule
\end{tabular}}
\end{center}

\noindent\footnotesize Latencies in milliseconds, end to end for Ours and call-plus-calibration
for the fixed rows (\S\ref{sec:lat}). \ours\,pool member.\normalsize

Ours is first of sixteen by AUROC and first by AURC ($.14905$ against $.15431$ for AlignScore).
It is separable from fourteen of the fifteen fixed verifiers. \fix{The exception is AlignScore:}
$+.01217$, nominally $[.00038,.02426]$, but the table makes fifteen comparisons and the
Bonferroni-corrected interval $[-.00721,.02964]$ contains zero. \fix{An earlier draft of these
notes claimed separability from all fifteen; that was wrong.} The honest statement is that Ours
holds the highest absolute AUROC and is indistinguishable from the strongest fixed verifier under
multiplicity correction, at 6.02$\times$ lower latency.

The three pool members individually reach $.53237$, $.61216$ and $.74024$; routing among them
reaches $.79330$.

\subsection{Protocol B}

\begin{center}\scriptsize
\setlength{\tabcolsep}{2.6pt}
\resizebox{\textwidth}{!}{%
\begin{tabular}{lrrrrrrrrrrrrr}
\toprule
System & AUROC$\uparrow$ & SD$\downarrow$ & AUPRC-U$\uparrow$ & Acc$\uparrow$ & BAcc$\uparrow$ & MCC$\uparrow$ & F1-U$\uparrow$ & Rec-U$\uparrow$ & ECE$\downarrow$ & Brier$\downarrow$ & Mean$\downarrow$ & P95$\downarrow$ & P99$\downarrow$ \\
\midrule
Qwen30-fast & \textbf{.83384} & .00066 & \textbf{.85299} & \textbf{.76681} & \textbf{.77251} & \textbf{.56311} & \textbf{.74098} & .64269 & \textbf{.02506} & \textbf{.15688} & 211.65 & 284.51 & 402.61 \\
\textbf{Ours} & .82256 & .00474 & .83522 & .73646 & .74153 & .50613 & .70888 & .62605 & .03883 & .16855 & \textbf{106.14} & 225.16 & 757.63 \\
Qwen30-judge & .81500 & .00158 & .81237 & .75034 & .75918 & .55993 & .69986 & .55767 & .05597 & .16377 & 600.87 & 1354.61 & 1426.40 \\
AlignScore & .81167 & .00108 & .83801 & .72951 & .73058 & .46608 & .73008 & \textbf{.70628} & .04345 & .17279 & 561.65 & 1747.54 & 5274.19 \\
Granite-3.1-2b\ours & .79583 & .00316 & .81944 & .72151 & .72866 & .48515 & .67718 & .56566 & .05570 & .18276 & 131.81 & 205.00 & 278.04 \\
FactCG & .77998 & .00080 & .82291 & .72886 & .73699 & .50752 & .67882 & .55169 & .09607 & .18862 & 427.50 & 1073.47 & 7859.26 \\
WeCheck & .76912 & .00098 & .80861 & .69385 & .70300 & .44893 & .62333 & .49449 & .09246 & .19609 & 496.01 & 1436.41 & 4536.65 \\
FactKB & .76282 & .00526 & .80264 & .70547 & .71733 & .50983 & .61214 & .44695 & .05945 & .18287 & \textbf{37.01} & \textbf{94.85} & 333.70 \\
Granite-3.2-3b & .75501 & .00792 & .75989 & .63276 & .64634 & .36790 & .48766 & .33683 & .14207 & .22169 & 263.87 & 456.21 & 549.55 \\
MiniCheck-FT5 & .73907 & .00279 & .76218 & .68847 & .69632 & .41990 & .63219 & .51741 & .12191 & .21586 & 364.03 & 1681.60 & 4196.31 \\
MiniCheck-DBTA & .67740 & .00135 & .69006 & .55732 & .57525 & .25710 & .28131 & .16679 & .12963 & .24017 & 480.64 & 2035.34 & 4264.31 \\
HHEM & .66521 & .01751 & .69026 & .60250 & .61900 & .35367 & .38750 & .24304 & .12600 & .23292 & 223.30 & 700.32 & 2428.56 \\
Granite-3.2-8b & .63386 & .00582 & .74705 & .53993 & .53511 & .06292 & .59823 & .64494 & .28108 & .26553 & 690.01 & 1176.09 & 1331.46 \\
Lettuce-v2\ours & .59951 & .00056 & .59493 & .58616 & .59992 & .25416 & .41920 & .28615 & .15271 & .25747 & 61.61 & 126.89 & 713.54 \\
FactCC\ours & .57089 & .00778 & .56859 & .53198 & .53392 & .09288 & .45340 & .48988 & .16127 & .27165 & 37.60 & 94.00 & 329.31 \\
Granite-4.1-3b & .47487 & .00179 & .53075 & .46378 & .46108 & $-$.07875 & .50382 & .52274 & .27710 & .33312 & 376.76 & 478.42 & 573.43 \\
\bottomrule
\end{tabular}}
\end{center}

Ours is second of sixteen. Paired intervals for Ours minus the arm, cluster bootstrap by source
group, 2{,}000 draws:

\begin{center}\small
\begin{tabular}{lrlll}
\toprule
Arm & $\Delta$ AUROC & 95\% CI & Bonferroni & Latency \\
\midrule
Granite-3.1-2b, strongest pool member & $+.02672$ & $[.00916,.04315]$ & \textbf{separable} & 1.24$\times$ faster \\
AlignScore & $+.01089$ & $[-.00625,.02753]$ & indistinguishable & 5.29$\times$ faster \\
Qwen30-judge & $+.00755$ & $[-.01044,.02477]$ & indistinguishable & 5.66$\times$ faster \\
Qwen30-fast & $-.01128$ & $[-.02875,.00552]$ & indistinguishable & 1.99$\times$ slower \\
\bottomrule
\end{tabular}
\end{center}

Twelve of the fifteen comparisons survive correction. The three that do not were already reported
as indistinguishable at the nominal level, so no conclusion reverses. In neither protocol is any
fixed verifier both more accurate and cheaper than Ours: under end-to-end accounting the only
cheaper verifiers are FactKB, FactCC and Lettuce-v2, all separably worse. The operating point is
not dominated. Universal superiority over every fixed verifier is not claimed and Qwen30-fast is
genuinely higher in Protocol~B.

\subsection{Routing behaviour}

\begin{center}\small
\begin{tabular}{lrrrr}
\toprule
Protocol & FactCC & Lettuce-v2 & Granite-3.1-2b & Mean ms \\
\midrule
A & 29.77\% & 51.32\% & 18.91\% & 112.54 \\
B & 34.50\% & 35.59\% & 29.91\% & 106.14 \\
\bottomrule
\end{tabular}
\end{center}

Every member carries at least 18.9\% of traffic in both protocols, so no pool slot is dead. The
per-corpus split is extreme --- in Protocol~B CogenSumm sends 88.7\% to FactCC, RAGTruth sends
98.9\% to Lettuce-v2, FRANK sends 55.0\% to Granite --- and is reported as a description of
fitted behaviour. It is not evidence of domain-invariant routing, and the cheap features carry a
strong corpus fingerprint.

\subsection{Pool provenance}
\label{sec:pool}

\fix{This subsection is new and it reports a negative result.} The pool was selected in the
earliest codebase on a matrix that shares 54.4\% of Protocol~B's TEST document groups. Features
and hyperparameters were later re-selected on TRAIN alone; the pool never was. Re-running the
selection on TRAIN only, over all 455 three-verifier subsets under the current features and
hyperparameters, and applying the project's own pre-declared eligibility rule ($k=3$, not
dominated by any single verifier, every member receiving at least 5\% of routed traffic):

\begin{center}\small
\begin{tabular}{lr}
\toprule
Validation AUROC & $.75930$ at $89.92$\,ms \\
Minimum member share & 13.6\%, eligible \\
Rank by validation AUROC alone & \fix{109 of 455} \\
Quality--cost Pareto frontier & \textbf{on it} (12 frontier pools of 455; 7 of the 90 eligible) \\
Seed SD among frontier pools & lowest \\
Cheapest pool with higher AUROC & 2.80$\times$ the cost \\
Best AUROC among cheaper pools & $.052$ lower \\
\bottomrule
\end{tabular}
\end{center}

The pool's rank-one status does not survive re-selection on clean data; its frontier position
does, with a wide margin on both axes. \fix{The previous claim of ``first of 540'' was measured on
the contaminated matrix and must not be quoted.}

%==================================================================
\section{Ablations}
\label{sec:abl}

Twenty-five single-factor arms, both protocols, ten seeds, each sharing the reference system's
splits and differing from it in exactly one respect. The reference arm reproduces the main
experiment at $\Delta=0$ to machine precision in both protocols, which is what licenses
attributing each delta to the change that produced it.

\subsection{What the method needs}

\begin{center}\small
\begin{tabular}{lrrl}
\toprule
Arm & B & A & \\
\midrule
Random routing, serving only & $-.16609$ & $-.14979$ & Bonferroni \\
Random routing, everything refitted on it & $-.16025$ & $-.14325$ & Bonferroni \\
Remove Granite at serving / retrained & $-.17629$ / $-.16425$ & $-.10244$ / $-.09893$ & Bonferroni \\
Remove Lettuce-v2 at serving / retrained & $-.08604$ / $-.02255$ & $-.09566$ / $-.05030$ & Bonferroni / 95\% \\
Remove FactCC at serving / retrained & $-.05377$ / $-.01840$ & $-.06703$ / $-.03843$ & Bonferroni \\
No verifier at all & $-.02589$ & $-.04405$ & Bonferroni \\
No stage-1 Platt & $-.04368$ & $-.03232$ & Bonferroni \\
\bottomrule
\end{tabular}
\end{center}

The routing decision is the dominant effect. The two random-routing variants matter jointly:
the serving-only version leaves beta, isotonic and the threshold fitted on real routing, so it
cannot separate the loss from a calibration mismatch; the refit version refits all three on
randomised routing and recovers only $.006$. The effect is the assignment of instances to
verifiers.

The two removal semantics answer different questions and are reported separately. Serving-time
masking is the production question of a verifier going down; retraining on the remainder is the
design question of whether a member earns its slot. Refitting recovers most of the loss for
FactCC and Lettuce-v2 --- they are partly substitutable --- and almost none for Granite.

\subsection{Three negative results}

\textbf{Stage-2 isotonic costs Protocol~B ranking quality.} Removing it \emph{improves} B's AUROC
by $+.00353$, separably under Bonferroni, while tripling its calibration error from $.03883$ to
$.11393$. In Protocol~A, where the stage also has to place ten rotations on a common scale, it
helps by $.00248$ at the 95\% level only. Stage 2 is a calibration component that costs a small,
measurable amount of ranking quality; it is not a component that improves the system.

\textbf{Individual feature attribution is weak.} Under Bonferroni only
\texttt{structured\_source\_line\_ratio} survives single-feature removal in Protocol~B
($-.03025$). Three of the six are not separable in either protocol, and two are nominally
positive when removed from B. The group arms keep the set defensible: dropping either
extractor's contribution is separable in both protocols, $-.03696$ and $-.02638$ in B. Section
\ref{sec:shap} resolves the attribution properly.

\textbf{The cost term's quality contribution is protocol-dependent.} Pinning $\beta=0$ costs
$-.03588$ in A, separable under Bonferroni, and $-.01352$ in B, separable at 95\% but not under
correction. It is a large latency regression in both, $+68\%$ and $+34\%$. The cost term is
primarily a budget control that also regularises over-selection of the expensive action.

\subsection{Supervision target}

\begin{center}\small
\begin{tabular}{lrrl}
\toprule
Target & B & A & \\
\midrule
\texttt{regret} (locked) & 0 & 0 & \\
\texttt{brier} & $-.00466$ & $-.00286$ & neither separable \\
\texttt{prob\_quality} & $-.00825$ & $-.00078$ & B Bonferroni; A not \\
\texttt{pairwise\_rank} & $-.01095$ & $+.00106$ & B Bonferroni; A not \\
\texttt{is\_best} & $-.03467$ & $-.01525$ & Bonferroni \\
\texttt{correctness} & $-.05699$ & $-.03118$ & Bonferroni \\
\bottomrule
\end{tabular}
\end{center}

Continuous relative-quality supervision beats binary supervision by $.015$ to $.057$, separably in
both protocols, which is the claim the target design supports. \fix{Optimality of \texttt{regret}
within the continuous family is not established:} it is separably best only in Protocol~B, and in
Protocol~A \texttt{pairwise\_rank} is nominally higher.

%==================================================================
\section{Extended ablation}

\subsection{Per-corpus training, and why the corpora are pooled}

Each corpus has its own official train and evaluation split. Training and evaluating entirely
inside one corpus, ten seeds, with the validation share raised until both classes and at least
five of the minority are present:

\begin{center}\small
\begin{tabular}{lrrrrr}
\toprule
Corpus & Train pos & Test pos & \textbf{Shift} & In-domain router & Pooled router, same rows \\
\midrule
CogenSumm & .4935 & .7800 & \textbf{$+28.65$ pp} & .68983 (sd .0399) & .68488 \\
RAGTruth & .7104 & .7822 & $+7.19$ pp & .62143 (sd .0165) & .63772 \\
UniSumEval & .7658 & .7466 & $-1.92$ pp & .63796 (sd .0281) & .59070 \\
FRANK & .1584 & .1638 & $+0.54$ pp & .81495 (sd .0117) & .82532 \\
\bottomrule
\end{tabular}
\end{center}

CogenSumm's own training and evaluation splits disagree on the label prior by 28.65 percentage
points. A model fitted inside that corpus cannot calibrate a threshold against itself. Pooling
reduces the worst-case shift a single fitted model must absorb to 15.17 points, and it is the
reason the four corpora are trained jointly rather than separately.

Single-corpus data buys nothing: the in-domain minus pooled differences are $+.005$, $-.016$,
$+.047$ and $-.010$, both signs, no systematic advantage, and the pooled router's seed SD is
$.0047$ against $.012$ to $.040$ in domain.

\subsubsection{A negative result that bounds the claim}

\fix{The deployable single-verifier competitor beats in-domain routing on every corpus.} Selecting
one fixed verifier on the same validation split, never touching test labels:

\begin{center}\small
\begin{tabular}{lrrrrr}
\toprule
Corpus & In-domain router & Selected from the 3 pool members & $\Delta$ & Selected from all 15 & $\Delta$ \\
\midrule
CogenSumm & .68983 & \textbf{.74948} Granite & $-.05965$ & .77583 WeCheck & $-.08600$ \\
FRANK & .81495 & \textbf{.82667} Granite & $-.01172$ & .83863 FactKB & $-.02367$ \\
RAGTruth & .62143 & \textbf{.63676} Lettuce & $-.01532$ & .78181 FactCG & $-.16038$ \\
UniSumEval & .63796 & \textbf{.66421} Granite & $-.02624$ & .66564 WeCheck & $-.02768$ \\
\bottomrule
\end{tabular}
\end{center}

Even restricted to the router's own three actions, fixing one verifier per corpus wins on all
four, and the selection is stable --- three of the four corpora choose the same verifier on all
ten seeds. Pooled, the ordering reverses: the router beats the strongest pool member by $+.02672$,
separably under Bonferroni, at 1.24$\times$ lower latency.

Two measurable conditions explain the reversal. Heterogeneity: four different verifiers win on
the four corpora when the choice is unrestricted (FactCG, FactKB, AlignScore, Qwen30-judge), so no
verifier is universally best across the mixture, which is what creates the arbitrage. Data: the
learning curve below is still rising at the full pooled training set of 890 groups, and a single
corpus has 107 to 499.

The defensible claim is therefore conditional. \textbf{Per-instance routing pays off on a
heterogeneous mixture, not within a single corpus. With one corpus and labels, fixing one verifier
for it is the stronger choice.}

\subsection{Feature attribution over the full subset lattice}
\label{sec:shap}

All 64 subsets of the six features and all 8 subsets of the pool were evaluated. With no features
the forest receives a single constant column, so each head predicts the training mean of its
action's target and routing reduces to the cost-discounted marginal ranking; that is the value of
the empty set which exact attribution requires.

\fix{Why this replaces per-subset significance testing.} Sixty-four comparisons would require a
Bonferroni threshold of $\alpha=.00078$ and would return ``not separable'' almost everywhere,
which is a property of the design rather than a finding, because the features are redundant. The
average marginal contribution over all subsets containing a feature is the quantity that answers
what a feature is worth under redundancy. It is exact at $n=6$, and the efficiency identity ---
the values sum to full-set minus empty-set --- holds to $0$ and is asserted at run time.

\begin{center}\small
\begin{tabular}{lrrrr}
\toprule
Feature & Shapley & Share & Alone & Leave-one-out \\
\midrule
\texttt{structured\_source\_line\_ratio} & \textbf{.01694} & \textbf{63.4\%} & .78785 & .79230 \\
\texttt{entity\_coverage} & .00267 & 10.0\% & .79382 & .82203 \\
\texttt{bm25\_mean3} & .00266 & 9.9\% & .76856 & .80767 \\
\texttt{conflicting\_value\_rate} & .00258 & 9.7\% & \textbf{.79587} & .82245 \\
\texttt{entity\_value\_colocation} & .00114 & 4.3\% & .78763 & .82175 \\
\texttt{year\_count} & .00074 & 2.8\% & .79465 & .82359 \\
\bottomrule
\end{tabular}
\end{center}

All six contributions are positive. The three that single-feature removal could not separate from
noise are small but not zero.

\textbf{Both simpler diagnostics mis-rank the features.} \texttt{conflicting\_value\_rate} is the
best single feature in isolation yet fourth by marginal contribution;
\texttt{structured\_source\_line\_ratio} is fifth in isolation yet accounts for 63.4\% of the
total. Under redundancy, neither solo performance nor leave-one-out recovers the ordering.

\begin{center}\small
\begin{tabular}{rrlrr}
\toprule
$k$ & Subsets & Best subset at that size & AUROC & Mean ms \\
\midrule
0 & 1 & no features & .79583 & 159.14 \\
1 & 6 & \texttt{conflicting\_value\_rate} & .79587 & 157.08 \\
2 & 15 & \texttt{structured\ldots} + \texttt{bm25\_mean3} & .81280 & 112.74 \\
3 & 20 & + \texttt{entity\_value\_colocation} & .82116 & 106.59 \\
4 & 15 & + \texttt{entity\_coverage} & .82317 & 105.25 \\
\textbf{5} & 6 & the full set minus \texttt{year\_count} & \textbf{.82359} & 106.68 \\
6 & 1 & the full set, as deployed & .82256 & 106.15 \\
\bottomrule
\end{tabular}
\end{center}

\fix{An honest negative result.} The best five-feature subset beats the deployed six-feature set
by $+.00103$, and the feature it drops is \texttt{year\_count}, matching the single-removal result
exactly. The forward-addition search that produced the current set overshot by one feature. The
margin is small and not separable, and the configuration is not changed after TEST was read, but
the fact is recorded.

The pool lattice gives a complementary observation. Shapley values are $+.12678$ for Granite,
$-.04223$ for Lettuce-v2 and $-.05862$ for FactCC, because the empty set --- predicting the label
from features with no verifier call --- already reaches $.79662$, while the best single verifier
reaches only $.79583$. \textbf{A single weak verifier is worse than no verifier at all, and the
pool's value is entirely in the combination.} The negative Shapley values and the positive
leave-one-out values are both correct: these two members are harmful in isolation and necessary in
the full pool.

\subsection{Convergence}
\label{sec:conv}

\begin{center}\small
\begin{tabular}{rrr@{\qquad}rrr}
\toprule
Fit groups & Val head loss & Eval AUROC & Trees & Val head loss & Eval AUROC \\
\midrule
5\% & .017293 & .790242 & 1 & .014154 & .806669 \\
10\% & .016430 & .790091 & 5 & .013502 & .816332 \\
20\% & .014115 & .806848 & 25 & .013341 & .822401 \\
35\% & .013695 & .808914 & \textbf{50} & \textbf{.013307} & \textbf{.825920} \\
50\% & .013849 & .813204 & 100 & .013302 & .825952 \\
75\% & .013117 & .817730 & 200 & .013290 & .826084 \\
100\% & .013290 & .822556 & & & \\
\bottomrule
\end{tabular}
\end{center}

\fix{Annotation volume is still the binding constraint}, contradicting the previous notes: the
curve is still rising at 100\%, gaining $.0048$ from 75\% to 100\%. The forests, by contrast, have
converged by about 50 trees; going from 50 to 200 gains $.000164$ AUROC. Protocol~B's 200 trees
are safely past convergence, and there is a fourfold saving available in head inference at no
measurable cost. The tree column is measured before stage 2, which is why its endpoint matches the
\texttt{no\_stage2} ablation rather than the main table.

\subsection{Why the feature sets were unified}
\label{sec:unify}

\begin{center}\small
\begin{tabular}{lrrll}
\toprule
Configuration & AUROC & $\Delta$ & 95\% CI & Bonferroni \\
\midrule
\textbf{unified six (deployed)} & \textbf{.82256} & 0 & --- & --- \\
superseded five-feature set & .81396 & $-.00860$ & $[.00214,.01489]$ & \textbf{separable} \\
superseded nine-feature set & .81053 & $-.01203$ & $[.00524,.01862]$ & \textbf{separable} \\
\bottomrule
\end{tabular}
\end{center}

Both superseded sets are separably worse under correction. Unification was a net gain, not a
simplification bought with quality.

\subsection{Two tight controls}

\begin{center}\small
\begin{tabular}{lrrl}
\toprule
Arm & AUROC & $\Delta$ & Mean ms \\
\midrule
reference & .82256 & 0 & 106.14 \\
call-proportion-matched random routing & .66074 & $-.16182$ & 102.85 \\
label oracle heads, $\beta^{*}$ unchanged & .99135 & $+.16880$ & 117.54 \\
label oracle heads, $\beta=0$ & .99421 & $+.17165$ & 127.78 \\
\bottomrule
\end{tabular}
\end{center}

The matched control samples actions from the router's own realised call proportions, so its budget
is the same to within $3.3$\,ms and the only thing destroyed is which instance goes to which
verifier. It loses $.162$.

The oracle bounds the headroom: with perfect head predictions the existing three verifiers reach
$.99135$. Of the $.169$ gap from the current system, \textbf{98.3\% is head prediction error and
1.7\% is the cost discount}. Relaxing the budget can buy back at most $.00285$. Future effort
belongs in the heads.

\subsection{[pending] Protocol A replication}

The Protocol~A counterparts of the lattice, the convergence curves, the unification contrast and
the controls were still computing when this was written. The two questions they settle are whether
the Shapley ordering is protocol-dependent, and whether the five-feature optimum reappears.

%==================================================================
\section{[pending] Cascades, second calls, and learner sensitivity}

Three families were running when this was written; Protocol~B is complete and its direction is
already unambiguous, so it is summarised here and will be tabulated when both protocols finish.

All three cascades --- escalate on the cheapest verifier's confidence, escalate when the two
cheapest disagree, escalate when a learned model predicts the cheapest will err --- are both
worse and more expensive than routing before calling, losing $.049$ to $.204$ AUROC while making
$1.53$ to $2.19$ calls per instance. Each enters through the weakest verifier, and a confidently
wrong first answer does not trigger escalation; the router declines to use that verifier before
paying for any score.

No second-call trigger produces a Bonferroni-separable gain, and each raises both calls and
latency, so the primary system keeps one call per instance. Seven alternative learners are all
separably worse than the random forest.

%==================================================================
\section{Limits and limitations}

\begin{enumerate}[label=\textbf{\arabic*.}]
\item \textbf{Multiplicity.} Each main table makes fifteen comparisons. Under correction
Protocol~A holds fourteen of fifteen and Protocol~B twelve of fifteen. No conclusion reverses,
but the corrected level is the one to cite.
\item \textbf{Pool provenance.} The pool was selected on a matrix sharing 54.4\% of Protocol~B's
test groups. It is on the Pareto frontier of a clean re-selection but ranks 109 of 455 on
validation AUROC alone (\S\ref{sec:pool}).
\item \textbf{Protocol A is cross-validated, not held out.} Its configuration was chosen on the
TRAIN subset of its own frame.
\item \textbf{Routing is not universally preferable.} Within a single corpus with labels, fixing
one verifier beats in-domain routing on all four corpora. The method's case is the heterogeneous
mixture.
\item \textbf{No transfer claim.} Neither protocol establishes transfer to an unseen corpus
family, and no leave-one-corpus-out experiment was run. The cheap features carry a strong corpus
fingerprint, so the pooled result must not be read as domain-invariant.
\item \textbf{Feature selection did not run the full pipeline.} Candidates were scored on
pre-stage-2 routed AUROC, and features were not re-selected after the hyperparameters changed.
The exhaustive lattice shows the selected set is one feature larger than the optimum.
\item \textbf{Two latency regimes.} Granite-3.1-2b's recorded latency differs by 1.5 to
1.7$\times$ between scoring campaigns. The conservative figure is reported (\S\ref{sec:lat}).
\item \textbf{Routing is not cheaper everywhere.} On UniSumEval the router averages $278.20$\,ms
against $147.60$ for always-Granite, because that corpus's long sources make FactCC's window
count grow.
\item \textbf{Action hit rate cannot decompose AUROC.} Agreement with the label oracle is a
row-wise three-class statistic; AUROC is determined by score ordering over pairs. No quality
decomposition is claimed from it.
\item \textbf{Per-instance $\beta$ is unproven.} $\beta$ is selected on validation from a finite
grid under a tolerance rule. That is an operational budget policy, not an optimality result.
\item \textbf{Macro AUROC is diagnostic only.} It is reported throughout and is never a ranking
key; the router is consistently weaker on it than several fixed verifiers.
\end{enumerate}

\end{document}
